\begin{afterword}
\summaryhead

The post opens by quoting the previous post's claim that one can never get Bayesian confirmation for a hypothesis of irreducibility. It then quotes Benja Fallenstein's objection: no test could logically prove irreducibility, but an outcome could make it much more probable under a simplicity prior, as a working hypercomputer might. The author concedes (``You're right'') and adds a similar flaw found independently.

If beliefs were ontologically basic, a direct law between minds would be simple. So telepathy, precognition and influencing events by wishing are a privileged probabilistic prediction of supernatural models. For reductionists, beliefs are complex patterns of neurons with known inputs and outputs, and nothing more is expected. A naive Occam rule that counts ``belief'' as one entity favours psychic powers; the naturalistic rule does not. A discovery of psychic powers would therefore be strong Bayesian evidence, though not proof, that beliefs are fundamental, and non-reductionists, who put the prediction forward, would win. The author concludes by dismissing psychic powers as a priori implausible, despite all the claimed experimental evidence.

\responsehead

The post opens with a concession, quoted accurately and credited by name, and it corrects the previous post at the point where that post went too far. Its main argument is sound in form. If one hypothesis makes an outcome much more likely than its rival does, observing the outcome favours it, even when the rival could also produce it. The history supports the claim that believers in non-physical minds put the prediction forward: a founder of the Society for Psychical Research hoped telepathy would show mind communicating with mind apart from known channels.

The problems are in the last sentence. It dismisses psychic powers ``despite all the claimed experimental evidence,'' and the post does not describe that evidence. The evidence existed and was contested. Meta-analyses of ganzfeld experiments reported an effect in 1994, found none in thirty further studies in 1999, and reported one again, pooling 79 studies, in 2001.

The sentence also links to ``I Defy the Data!'' for its confidence. That post's addendum allows defying the data of one experiment, not of many. Past that point, it says, one must give up the theory or dismiss the data with reasons, such as a design flaw or a larger body of failed replications. Psychic powers rest on many claimed experiments, and the post gives neither kind of reason, although a null result of the kind the addendum names was in print.

In short: a prompt, credited concession and a sound argument that psychic powers would favour non-reductionism, ending in a dismissal that sets aside many claimed experiments without describing them, and that invokes a rule about defying data whose own addendum does not cover many experiments.
\end{afterword}
